\begin{center}
\hypertarget{native-047-article032}{}\pdfbookmark[2]{32.}{article-032}32.
\end{center}
Progredimur iam ad applicationem horum principiorum ad utilissimas operationum supra explicatarum.

I. Adhibendo ad computum anomaliae verae ex anomalia excentrica in motu elliptico formulam VII art. 8, si $\varphi$ et $E$ exacte haberi supponuntur, in \SourceLeaf{048}$\log\operatorname{tang}(45^0-\frac{1}{2}\varphi)$ et $\log\operatorname{tang}\frac{1}{2}E$ committi potest error $\omega$, adeoque in differentia $=\log\operatorname{tang}\frac{1}{2}v$ error $2\omega$; error maximus itaque in determinatione anguli $\frac{1}{2}v$ erit $\frac{3\omega\,\mathrm{d}\frac{1}{2}v}{\mathrm{d}\log\operatorname{tang}\frac{1}{2}v}=\frac{3\omega\sin v}{2\lambda}$, designante $\lambda$ modulum logarithmorum ad hunc calculum adhibitorum. Error itaque, cui anomalia vera $v$ obnoxia est, in secundis expressus fit $=\frac{3\omega\sin v}{\lambda}206265^{\prime\prime}=0^{\prime\prime}{,}0712\sin v$, si logarithmi \textsc{Brigg}ici ad septem figuras decimales adhibentur, ita ut semper intra $0^{\prime\prime}{,}07$ de valore ipsius $v$ certi esse possimus: si tabulae minores ad quinque tantum figuras adhibentur, error usque ad $7^{\prime\prime}{,}12$ ascendere posset.

II. Si $e\cos E$ adiumento logarithmorum computatur, error committi potest usque ad $\frac{3\omega e\cos E}{\lambda}$; eidem itaque errori obnoxia erit quantitas $1-e\cos E$ sive $\frac{r}{a}$. In computando ergo logarithmo huius quantitatis error usque ad $(1+\delta)\omega$ ascendere potest, designando per $\delta$ quantitatem $\frac{3e\cos E}{1-e\cos E}$ positive sumtam: ad eundem limitem $(1+\delta)\omega$ ascendit error in $\log r$ possibilis, siquidem $\log a$ exacte datus supponitur. Quoties excentricitas parva est, quantitas $\delta$ arctis semper limitibus coërcetur: quando vero $e$ parum differt ab 1, $1-e\cos E$ perparva manet, quamdiu $E$ parva est; tunc igitur $\delta$ ad magnitudinem haud contemnendam increscere potest, quocirca in hoc casu formula III art. 8 minus idonea esset. Quantitas $\delta$ ita etiam exprimi potest $\frac{3(a-r)}{r}=\frac{3e(\cos v+e)}{1-ee}$, quae formula adhuc clarius ostendit, quando errorem $(1+\delta)\omega$ contemnere liceat.

III. Adhibendo formulam X art. 8 ad computum anomaliae verae ex excentrica, $\log\BareRoot{\frac{a}{r}}$ obnoxius erit errori $(\frac{1}{2}+\frac{1}{2}\delta)\omega$, adeoque $\log\sin\frac{1}{2}\varphi\sin E\BareRoot{\frac{a}{r}}$ huic $(\frac{5}{2}+\frac{1}{2}\delta)\omega$; hinc error maximus in determinatione anguli $v-E$ vel $v$ possibilis eruitur $=\frac{\omega}{\lambda}(7+\delta)\operatorname{tang}\frac{1}{2}(v-E)$, sive in secundis expressus, si septem figurae decimales adhibentur, $=(0^{\prime\prime}{,}166+0^{\prime\prime}{,}024\,\delta)\operatorname{tang}\frac{1}{2}(v-E)$. Quoties excentricitas modica est, $\delta$ et $\operatorname{tang}\frac{1}{2}(v-E)$ quantitates parvae erunt, quapropter haec methodus praecisionem maiorem permittet, quam ea quam in I contemplati sumus: haecce contra methodus tunc praeferenda erit, quando excentricitas valde magna est propeque ad unitatem accedit, ubi $\delta$ et $\operatorname{tang}\frac{1}{2}(v-E)$ valores valde considerabiles nancisci possunt. Per formulas nostras, utra methodus alteri praeferenda sit, facile semper decidi poterit.

IV. In determinatione anomaliae mediae ex excentrica per formulam XII art. 8 error quantitatis $e\sin E$, adiumento logarithmorum computatae, adeoque etiam ipsius anomaliae $M$, usque ad $\frac{3\omega e\sin E}{\lambda}$ ascendere potest, qui erroris limes
\SourceLeaf{049}si in secundis expressus desideratur per $206265^{\prime\prime}$ est multiplicandus. Hinc facile concluditur, in problemate inverso, ubi $E$ ex $M$ tentando determinatur, $E$ quantitate $\frac{3\omega e\sin E}{\lambda}\cdot\frac{\mathrm{d}E}{\mathrm{d}M}\cdot206265^{\prime\prime}=\frac{3\omega ea\sin E}{\lambda r}\cdot206265^{\prime\prime}$ erroneam esse posse, etsi aequationi $E-e\sin E=M$ omni quam tabulae permittunt praecisione satisfactum fuerit.

Anomalia vera itaque e media computata duabus rationibus erronea esse potest, siquidem mediam tamquam exacte datam consideramus, primo propter errorem in computo ipsius $v$ ex $E$ commissum, qui ut vidimus levis semper momenti est, secundo ideo quod valor anomaliae excentricae ipse iam erroneus esse potuit. Effectus rationis posterioris definietur per productum erroris in $E$ commissi per $\frac{\mathrm{d}v}{\mathrm{d}E}$, quod productum fit $=\frac{3\omega e\sin E}{\lambda}\cdot\frac{\mathrm{d}v}{\mathrm{d}M}\cdot206265^{\prime\prime}=\frac{3\omega ea\sin v}{\lambda r}\cdot206265^{\prime\prime}=\left(\frac{e\sin v+\frac{1}{2}ee\sin 2v}{1-ee}\right)0^{\prime\prime}{,}0712$, si septem figurae adhibentur. Hic error, pro valoribus parvis ipsius $e$ semper modicus, permagnus evadere potest, quoties $e$ ab unitate parum differt, uti tabella sequens ostendit, quae pro quibusdam valoribus ipsius $e$ valorem maximum illius expressionis exhibet.

\begin{center}
\begin{tabular}{c|c||c|c||c|c}
$e$ & error maximus & $e$ & error maximus & $e$ & error maximus \\
0,90 & $0^{\prime\prime}{,}42$ & 0,94 & $0^{\prime\prime}{,}73$ & 0,98 & $2^{\prime\prime}{,}28$ \\
0,91 & $0{,}48$ & 0,95 & $0{,}89$ & 0,99 & $4{,}59$ \\
0,92 & $0{,}54$ & 0,96 & $1{,}12$ & 0,999 & $46{,}23$ \\
0,93 & $0{,}62$ & 0,97 & $1{,}50$ & & \\
\end{tabular}
\end{center}

V. In motu hyperbolico, si $v$ per formulam III art. 21 ex $F$ et $\psi$ exacte notis determinatur, error usque ad $\frac{3\omega\sin v}{\lambda}\cdot206265^{\prime\prime}$ ascendere potest; si vero per formulam $\operatorname{tang}\frac{1}{2}v=\frac{(u-1)\operatorname{cotang}\frac{1}{2}\psi}{u+1}$ computatur, $u$ et $\psi$ exacte notis, erroris limes triente maior erit, puta $=\frac{4\omega\sin v}{\lambda}\cdot206265^{\prime\prime}=0^{\prime\prime}{,}09\sin v$ pro septem figuris.

VI. Si per formulam XI art. 22 quantitas $\frac{\lambda kt}{b^{\frac{3}{2}}}=N$ adiumento logarithmorum \textsc{Brigg}icorum computatur, $e$ et $u$ vel $e$ et $F$ tamquam exacte notas supponendo, pars prima obnoxia erit errori $\frac{5(uu-1)e\omega}{2u}$, si computata est in forma $\frac{\lambda e(u-1)(u+1)}{2u}$, vel errori $\frac{3(uu+1)e\omega}{2u}$, si computata est in forma $\frac{1}{2}\lambda eu-\frac{\lambda e}{2u}$, vel errori $3e\omega\operatorname{tang}F$, si computata est in forma $\lambda e\operatorname{tang}F$, siquidem errorem in $\log\lambda$ vel $\log\frac{1}{2}\lambda$ commissum contemnimus. In casu primo error etiam per
\SourceLeaf{050}$5e\omega\operatorname{tang}F$, in secundo per $\frac{3e\omega}{\cos F}$ exprimi potest, unde patet, in casu tertio errorem omnium semper minimum esse, in primo autem vel secundo maior erit, prout $u$ aut $\frac{1}{u}>2$ vel $<2$, sive prout $\pm F>36^0 52^{\prime}$ vel $<36^0 52^{\prime}$. — Pars secunda ipsius $N$ autem semper obnoxia erit errori $\omega$.

VII. Vice versa patet, si $u$ vel $F$ ex $N$ tentando eruatur, $u$ obnoxiam fore errori $(1\pm5e\operatorname{tang}F)\omega\frac{\mathrm{d}u}{\mathrm{d}N}$, vel huic $\left(1+\frac{3e}{\cos F}\right)\omega\frac{\mathrm{d}u}{\mathrm{d}N}$, prout membrum primum in valore ipsius $N$ vel in factores vel in partes resolutum adhibeatur; $F$ autem errori huic $(1\pm3e\operatorname{tang}F)\omega\frac{\mathrm{d}F}{\mathrm{d}N}$. Signa superiora post perihelium, inferiora ante perihelium valent. Quodsi hic pro $\frac{\mathrm{d}u}{\mathrm{d}N}$ vel pro $\frac{\mathrm{d}F}{\mathrm{d}N}$ substituitur $\frac{\mathrm{d}v}{\mathrm{d}N}$, emerget effectus huius erroris in determinationem ipsius $v$, qui igitur erit $\frac{bb\operatorname{tang}\psi(1\pm5e\operatorname{tang}F)\omega}{\lambda rr}$ aut $\frac{bb\operatorname{tang}\psi(1+3e\operatorname{sec} F)\omega}{\lambda rr}$, si quantitas auxiliaris $u$ adhibita est; contra, si adhibita est $F$, ille effectus fit
\[
\begin{aligned}
&=\frac{bb\operatorname{tang}\psi(1\pm3e\operatorname{tang}F)\omega}{\lambda rr}\\
&=\frac{\omega}{\lambda}\left\{\frac{(1+e\cos v)^2}{\operatorname{tang}\psi^3}\pm\frac{3e\sin v(1+e\cos v)}{\operatorname{tang}\psi^2}\right\}\raisebox{.45ex}{.}
\end{aligned}
\]
Adiicere oportet factorem $206265^{\prime\prime}$, si error in secundis exprimendus est. Manifesto hic error tunc tantum considerabilis evadere potest, quando $\psi$ est angulus parvus, sive $e$ paullo maior quam 1; ecce valores maximos huius tertiae expressionis pro quibusdam valoribus ipsius $e$, si septem figurae decimales adhibentur:

\begin{center}
\begin{tabular}{c|c}
$e$ & error maximus \\
1,3 & $0^{\prime\prime}{,}34$ \\
1,2 & $0{,}54$ \\
1,1 & $1{,}31$ \\
1,05 & $3{,}36$ \\
1,01 & $34{,}41$ \\
1,001 & $1064{,}65$ \\
\end{tabular}
\end{center}
\noindent Huic errori ex erroneo valore ipsius $F$ vel $u$ orto adiicere oportet errorem in V determinatum, ut incertitudo totalis ipsius $v$ habeatur.

VIII. Si aequatio XI art. 22 adiumento logarithmorum hyperbolicorum solvitur, $F$ pro quantitate auxiliari adhibita, effectus erroris in hac operatione possibilis in determinationem ipsius $v$ per similia ratiocinia invenitur
\SourceLeaf{051}
\[
= \frac{(1+e\cos v)^2\,\omega'}{\mathrm{tang}\,\psi^3} \pm \frac{3e\sin v(1+e\cos v)\,\omega}{\lambda\,\mathrm{tang}\,\psi^2},
\]
ubi per $\omega'$ incertitudinem maximam in tabulis logarithmorum hyperbolicorum designamus. Pars secunda huius expressionis identica est cum parte secunda expressionis in VII traditae, prima vero in ratione $\lambda\omega' : \omega$ minor quam prima in illa expressione, i.\,e. in ratione $1:23$, si tabulam \textsc{Ursin}i ad octo ubique figuras exactam sive $\omega' = 0{,}000\,0000\,05$ supponere liceret.
